% =============================================================================
% Coda -- The map in both directions
% =============================================================================
\noindent The construction is complete. It is worth saying, in plain language and
for the reader who came to this document from philosophy of physics rather than
from lattice field theory, what the finished map shows, and what it honestly
does not.

\runin{The two findings, stated together.} A finite, discrete ontology\,---\,five
postulates, no completed infinity, no primitive continuum, together with the
single $\Zi$ reading of its order-4 symmetry (\secref{sec:seed})\,---\,\textbf{forces}
a rich theorem-grade algebraic spine: the constant $\Gs=\Gamma(1/4)/\Gamma(3/4)$
(in four representations of one CM period), the master quadratic and its
roots, the arithmetic of the lemniscatic CM curve, the harmonic-invariant tower,
the $\pi$-freeness of $\Q(\Gs)$: \textbf{seven theorem-grade results} in all
(Part~I). And the \emph{same} ontology, in the precise conditional sense of
Part~II, \textbf{does not select} the electromagnetic coupling $\alpha$ without an
additional binding law: the boundary is route-invariant, built from
theorem-grade masonry around one precisely-located open kernel, and its honest
verdict is that $\alpha$ is \emph{dynamical, not structural}. \textbf{Both
halves are the deliverable.} This is the point that is easiest to miss and most
important to hold: the second finding is not the first finding's failure. A
mapped boundary (\emph{this} is what the ontology reaches, and \emph{this} is the
precise place it stops under the present axioms) is a result in its own right, a
different and non-trivial kind of result, one easy to mistake for a failure. A
program that exhibited only the spine would be a curiosity; one that maps the
wall, in both directions and at honest tags, is doing the thing the goal was set
to do.

\runin{The least-wrong self-assessment.} Stated without flattery: FTD is a
\textbf{philosophy-of-mathematics project with a rigorous algebraic core and
suggestive (not derived) physics connections.} The algebraic core is
genuine: the seven theorem-grade results survive skeptical mathematical review,
are independently checkable, and stand independent of any physical
interpretation. The physics connection is real and is honestly weaker than the
mathematics: the $\alpha$ match is a \etag{smc} backed by strong
structural-uniqueness evidence and no derivation chain; the broader mass-and-gauge
layer is \etag{par}. FTD is therefore \textbf{not} a competitor to QED or general
relativity on their empirical terms; it does not out-predict them, and it
does not claim to. Its contribution is of a different kind: a demonstration,
carried out at theorem grade where it can be and at honest conjecture grade where
it cannot, of exactly how far a discrete ontology reaches toward physics on its
own, and exactly where it needs an additional input it does not contain.

\runin{Where it touches the open questions.} The program brushes against two
genuinely open foundational questions, and it touches each at that question's
honest status, claiming no resolution. On \textbf{the measurement problem}, FTD's
quantum-mechanics-emergence layer (the Born rule, collapse-as-projection) is
\etag{sel}/\etag{opn}, not a solution: the load-bearing step ``probability
$=$ normalized energy density'' is unproven, and the framework records it as
unproven. On \textbf{the status of mathematical structure in physics} (whether
the unreasonable effectiveness of mathematics reflects something forced or
something chosen) FTD offers an unusually crisp data point precisely
\emph{because} it draws the line: here is a case where one can exhibit, with
proofs, that a discrete mathematical ontology forces a specific rich structure,
and can exhibit, with theorem-grade conditional machinery, where the same
ontology stops short of a specific physical coupling. The effectiveness is real
on the forced side and absent on the unforced side, and the boundary between them
is mapped rather than asserted. That is the most the project claims, and it
claims it at the tags the claim can bear.

\runin{The two-clause goal, closed.} The value of this construction is not a
derivation of physics from discreteness; that, Part~II shows, the ontology
does not supply. The value is the \textbf{map}: a faithful drawing of the line of
what a finite discrete ontology reaches, \textbf{honest in both
directions}: every theorem proved where the ontology forces, every boundary
marked where it does not, and no tag promoted to make the picture prettier than
it is. The constant at the centre of it all, $\Gs\approx\Gstarapprox$, is the
lemniscatic ratio $\Gamma(1/4)/\Gamma(3/4)$ throughout, never the distinct
lemniscate constant $\varpi\approx\varpiapprox$. That is the construction, and
that is its honest reach.
